\begin{abstract}
  标准 FashionMNIST 的服饰主体大多位于图像中央，常规测试准确率因而难以反映模型面对位置变化时的稳定性。本项目将原始 $28\times28$ 灰度图像放入 $56\times56$ 画布，并施加可控平移、最大 $45^\circ$ 的双向旋转和随机擦除。由此构造的 PV-FashionMNIST 支持中心、小幅平移、大幅平移、$7\times7$ 位置网格、随机旋转和固定角度扫描等评估协议。在统一的 CPU 训练预算下，项目比较了 MLP、CNN、带可学习绝对位置编码的 Tiny ViT，以及依次加入复合增强、Mean Pooling 和卷积 stem 的三种 ViT 变体。

  实验显示，中心准确率不能替代位置鲁棒性评价。CNN 的 Center Accuracy 达到 $90.79\%$，但 Grid Accuracy 仅为 $13.36\%$；ViT-AbsPos 的 Grid Accuracy 为 $22.14\%$，虽高于另外两种中心训练基线，仍存在明显的中心高值区。扩大训练位置与角度分布后，ViT-Aug 和 ViT-MeanPool 的网格表现分别提升至 $62.29\%$ 和 $66.20\%$。加入卷积 stem 的 HybridConv-ViT 取得 $80.99\%$ 的 Grid Accuracy，49 个位置的准确率落在 $78.15\%$--$82.65\%$ 之间；其七个固定角度的平均准确率为 $80.53\%$。位置热力图、角度扇形图、逐类准确率、混淆矩阵和错误样本进一步说明：训练分布覆盖是缓解中心偏置的首要因素，局部卷积特征有助于提升小图像任务中的整体稳定性，但边缘裁剪和相似服饰类别仍构成主要困难。项目最终形成了一套从数据构造、模型训练、鲁棒性评估到可视化分析的可复现工程流程。\\

  \noindent\textbf{关键词：}位置可变性；FashionMNIST；Vision Transformer；鲁棒性评估；数据增强；可视化分析
\end{abstract}
